La masse d'un échantillon de soufre $S$ est $m=\qty{8}{\g}$.
\begin{enumerate}
    \item Calculer la quantité de matière qui se trouve dans cette
          échantillon.
    \item Déterminer le nombre d'atome qui se trouve dans cette masse.
\end{enumerate}

\begin{donnees}
$M(S)=\qty{32}{\g\per\mol}$, \ $N_A=\qty{6.02e23}{\per\mol}$
\end{donnees}
